\documentclass{article}
\usepackage{../../math_template}

\author{Jonathan Kasongo}
\title{Noether Sheet 2}

\begin{document}
\maketitle
\centerline{
These were my solutions to the Noether Mentoring Scheme Sheet 2. Problem
statements are abridged.
}

\begin{problem}{Problem 1}
Let $\Gamma$ and $\Delta$ be two circles that intersect in two points,
$P$ and $Q$. Suppose that $ST$ is tangent to $\Gamma$ at $S$ and
$\Delta$ at $T$. Prove that the extension of $PQ$ bisects $ST$.
\end{problem}

\begin{solution}{Solution}
\begin{center}
\begin{asy}
/* Geogebra to Asymptote conversion, documentation at artofproblemsolving.com/Wiki go to User:Azjps/geogebra */
import graph; size(150);
real labelscalefactor = 1.0; /* changes label-to-point distance */
real xmin = -8.020421870644345, xmax = 9.993974178973335, ymin = -5.700867645591302, ymax = 6.76314691846839;  /* image dimensions */
pen qqzzff = rgb(0.,0.6,1.); pen ttttff = rgb(0.2,0.2,1.); pen ffqqtt = rgb(1.,0.,0.2);
pair Q = (-1.1403008965047974,-1.4278995149108797), P = (-1.1403008965047974,1.42789951491088), T = (0.47700789900752283,3.0952999540571318), K = (-1.1403008965047974,2.289811351189327);

filldraw(circle((-2.,0.), 1.6667272042038497), qqzzff + opacity(0.15000000596046448), qqzzff);
filldraw(circle((2.,0.), 3.4496937176032834), ttttff + opacity(0.20000000298023224), ttttff);
/* draw figures */
draw((-2.7137189873216903,1.5061821869779506)--T, ffqqtt);
draw(Q--P,  ffqqtt);
draw(K--P, ffqqtt);
/* dots and labels */
label("$Q$", (-0.9815922172802164,-1.5137377529774978), NE * labelscalefactor);
label("$P$", (-0.9955029478204617,1.4631585826349733), NE * labelscalefactor);
//dot((-2.7137189873216903,1.5061821869779506),dotstyle);
label("$S$", (-2.6647906126498992,1.6439980796581608), NE * labelscalefactor);
label("$T$", (0.5346774116065228,3.229821361246113), NE * labelscalefactor);
label("$K$", (-1.0789673310619337,2.3951775288314012), NE * labelscalefactor);
clip((xmin,ymin)--(xmin,ymax)--(xmax,ymax)--(xmax,ymin)--cycle);
/* end of picture */
\end{asy}
\end{center}

Let $PQ$ and $ST$ intersect at the point $K$.
By Power of a point, $SK^2 = QK\cdot PK$ and $TK^2 = QK\cdot PK$. Since
both $SK>0$ and $TK>0$ we can conclude that $TK = SK$, and so the extension
of $PQ$ does in fact bisect $ST$. $\blacksquare$ \\

\textit{Comments: This is just a straight-forward application of
Power of a point, I had definitely heard of the theorem in the past but
now I know how to use it properly!}
\end{solution}
\newpage

\begin{problem}{Problem 2}
\begin{enumerate}[label=(\alph*)]
\item Show that $x^3 + y^3 = (x+y)(x^2-xy+y^2)$ and that
$x^3 - y^3 = (x-y)(x^2+xy+y^2)$
\item Find the prime factorisation of 1343.
\item Find the prime factorisation of 7657.
\end{enumerate}
\end{problem}

\begin{solution}{Solution}
\textbf{(a)} Expand.
$$(x+y)(x^2-xy+y^2) = (x^3 - x^2y +xy^2)+(yx^2-xy^2+y^3) = x^3 + y^3$$
$$(x-y)(x^2+xy+y^2) = (x^3 + x^2y + xy^2)-(yx^2+xy^2+y^3) = x^3 - y^3.
\ \square$$
\\
\textbf{(b)} Notice that $1343 = 10^3 + 7^3$ so we can use that identity
to write $1343 = (10+7)\cdot (10^2 - 7\cdot 10 + 7^2) = 17\cdot 79$. Now
one can check all divisors from 2 to 8 to verify whether or not these
numbers are prime. This is because to test whether a number $n$
is prime we need only to check numbers up to $\floor{\sqrt{n}}$ because
if a divisor $d$ of $n$ is greater than $\sqrt{n}$ then there exists another
divisor of $n$ that is less than $\sqrt{n}$. Indeed
$\frac{n}{d} < \sqrt{n} \iff \sqrt{n} < d$ which is true
and so $\frac{n}{d}$ is
a divisor of $n$ that is $< \sqrt{n}$.
\vspace{0.1cm}
\begin{center}
\begin{tblr}{
  colspec = {c c c},
  vlines,
}
$n$ & $17 \bmod n$ & $79 \bmod n$ \\
\hline
2 & 1 & 1 \\
3 & 2 & 1 \\
4 & 1 & 3 \\
5 & 2 & 4 \\
6 & 5 & 1 \\
7 & 3 & 2 \\
8 & 1 & 7 \\
\end{tblr}
\end{center}
and thus both 17 and 79 are prime so the prime factorisation of 1343 is
$17\cdot 79$. $\square$\\

\textbf{(c)} Notice that $7657 = 20^3 - 7^3 = (20-7)(20^2 + 20\cdot7 + 7^2)
= 13 \cdot 589 = 13 \cdot 19 \cdot 31$. This time we only need to check
divisors from 2 to 5. \vspace{0.1cm}

\begin{center}
\begin{tblr}{
  colspec = {c c c c},
  vlines,
}
$n$ & $13 \bmod n$ & $19 \bmod n$ & $31 \bmod n$ \\
\hline
2 & 1 & 1 & 1 \\
3 & 1 & 1 & 1 \\
4 & 1 & 3 & 3 \\
5 & 3 & 4 & 1 \\
\end{tblr}
\end{center}
so we have verified that 13,19 and 31 are all prime and so the desired
prime factorisation is $13\cdot 19\cdot 31$. $\blacksquare$
\end{solution}

\newpage

\begin{problem}{Problem 3}
\begin{enumerate}[label=(\alph*)]
  \item Derive the formula for $\binom{n}{r}$.
  \item A frog is hopping along a line. Every minute, it hops either 10 cm
  to the right or 10 cm to the left. After 10 minutes, the frog is back in
  the same position as it started. In how many ways could this
  be achieved?
  \item A celebrity has 10 identical signed photographs to distribute
  amongst five fans. (As the photographs are identical, the only thing that
  matters is the number of photographs that each fan gets.)
  \begin{enumerate}[label=\roman*.]
    \item How many different ways are there of doing this if the celebrity
    may decide to give no photographs to some of these fans?
    \item How many different ways are there of doing this if the celebrity
    gives at least one photograph to each fan?
  \end{enumerate}
\end{enumerate}
\end{problem}

\begin{solution}{Solution}
\textbf{(a)} Let $A$ be a set of size $n$. Recall the definition for
$\binom{n}{r}$, it is the number of ways to choose a subset
$B \subset A$ such that $|B| = r$. It is important to remember that
sets are unordered and that they don't contain repititions, by definition.
\vspace{0.2cm}

Notice that after selecting $k$ objects there remains $n-k$ objects to
choose from, so the number of ways to choose an ordered $r$-tuple of
distinct elements from $A$ without repitition is
$n(n-1)\cdots(n-r+1)$. Combinations are the number of ways to choose an
unordered $r$-tuple of distinct elements from $A$
without repitition, and since each ordered $r$-tuple has $r!$ ways to
permute we need to divide by $r!$. Thus

$$
\binom{n}{r} = \frac{n(n-1)\cdots(n-r+1)}{r!} = \frac{n!}{r!(n-r)!}. \
\square
$$
\\
\textbf{(b)} Suppose the frog hops left $l$ times and hops right $r$ times.
Let the frog hop along the $x$-axis in the Cartesian plane such
that 10 cm corresponds to a unit of length 1. If we let the frog start on
the origin it's final position will be at $(r-l,0)$. Since the
frog returns back to the origin it is the case that $r-l=0$ and so $r=l$.
\vspace{0.2cm}

Thus the frog jumps left 5 times, and jumps right 5 times. Let $L$
represent a move to the left, and let $R$ represent a move to the right.
We want to count the number of ordered $10$-tuples that consist of
exactly 5 $L$'s and exactly 5 $R$'s. There are $\binom{10}{5}$ ways to
place the $R$'s and the placement of the $R$'s uniquely determines the
placement of the $L$'s so the final answer is $\binom{10}{5} = 252$.
$\square$\\

\textbf{(c) i.} Let's represent the 5 fans by
$f_1, f_2, f_3, f_4, f_5 \in \mathbb{Z}_{\geq 0}$ where
the $i$-th fan recieves $f_i$ photos. We observe that
$$
f_1 + f_2 + f_3 + f_4 + f_5 = 10.
$$

and so we can use Stars and bars to show that the answer is
$\binom{14}{4}$. Imagine placing 10 stars and 4 bars in a line in some
order, for instance
$$
\text{
\Large
$\mid$
$\star$
$\star$
$\star$
$\mid$
$\star$
$\star$
$\star$
$\mid$
$\mid$
$\star$
$\star$
$\star$
$\star$
}
$$
\normalsize

The 4 bars separate the stars into 5 groups of $n\in \mathbb{Z}_{\geq 0}$
stars and we can let each group
represent a value of $f_i$. As an example the configuration above would
represent $(f_1, f_2, \ldots, f_5)=(0,3,3,0,4)$. Now since there are
$\binom{14}{4}$ ways to arrange the stars and bars our answer is
$\binom{14}{4} = 1001$. $\square$ \\

\textbf{(c) ii.} The solution for ii is similar to i. We use stars and
bars, except this time it is not permissable to have two bars next to
each other (as this would represent a group of size 0), neither can we have
any bars at one of the two endpoints of the list (this would also create
a group of size 0). We can instead write out all of the 10 stars in a row,
and in the gap between each star there may or may not be exactly one bar.
Here's an example

$$
\text{
\Large
$\star$
$\star$
$\mid$
$\star$
$\mid$
$\star$
$\mid$
$\star$
$\star$
$\star$
$\mid$
$\star$
$\star$
$\star$
}
$$

The 4 bars separate the stars into 5 groups of
$n \in \mathbb{N}$ stars and so we can let each group represent a value of
$f_i$.
and this configuration would represent $(f_1,f_2,\ldots, f_5) =
(2,1,1,3,3)$. There are 9 positions in which we could put the bars and 4
total bars to place so the answer is $\binom{9}{4} = 126$. $\blacksquare$
\\

\textit{Comments: Recently I have been studying combinatorics in
preparation for the TMUA, so I've had practise on these kinds of problems.}
\end{solution}
\newpage


\begin{problem}{Problem 4}
A triangle has side lengths $a, b$ and $c$.
In terms of these variables, what is the height of the triangle
(taking the base to be the side with length $c$)? \vspace{0.2cm}

Hence derive Heron's formula.
\end{problem}

\begin{solution}{Solution}
Consider,

\begin{center}
\begin{asy}
/* Geogebra to Asymptote conversion, documentation at artofproblemsolving.com/Wiki go to User:Azjps/geogebra */
import graph; size(140);
real labelscalefactor = 1.0; /* changes label-to-point distance */
pen dotstyle = black; /* point style */
real xmin = -4.374399355946519, xmax = 3.8826150852061123, ymin = -1.1598067106128818, ymax = 4.553154632454909;  /* image dimensions */
pen zzwwff = rgb(0.8,0.2,0.1); pen qqccqq = rgb(0.2,0.,0.9);
pair A = (-2.998686686773226,0.), B = (3.,0.), C = (-1.,4.), D = (-0.9991244578488174,0.);

filldraw(A--C--B--cycle, qqccqq + opacity(0.10000000149011612), qqccqq);
/* draw figures */
draw(C--A);
draw(C--B);
draw(A--B);
draw(C--D, linetype("4 4") + zzwwff);
label("$x$",(-2.0598848832527707,0.032518926835418224),SE*labelscalefactor);
label("$c-x$",(1.0197583407446964,0.05802321854554229),SE*labelscalefactor);
draw(A--C, qqccqq);
draw(C--B, qqccqq);
draw(B--A, qqccqq);
/* dots and labels */
label("$A$", (-3.531065135077977,0.05164714561801127), NE * labelscalefactor);
label("$B$", (2.85,0.06439929147307331), NE * labelscalefactor);
label("$C$", (-0.9759524855725027,4.062197017035021), NE * labelscalefactor);
label("$b$", (-2.5917652478904256,2.0473579719352193), NE * labelscalefactor);
label("$a$", (0.9241172468317317,1.932588659239661), NE * labelscalefactor);
label("$c$", (-4.1332766026157513e-4,-0.10137860464273311), NE * labelscalefactor);
label("$D$", (-0.9759524855725027,0.06439929147307331), NE * labelscalefactor);
label("$h$", (-1.1034739441231225,2.0027254614425023), NE * labelscalefactor,zzwwff);
clip((xmin,ymin)--(xmin,ymax)--(xmax,ymax)--(xmax,ymin)--cycle);
/* end of picture */
\end{asy}
\end{center}

due to Pythagoras theorem we have
\[
\begin{aligned}
b^2 - x^2 &= a^2 - (c-x)^2 \\
b^2 &= a^2 - c^2 + 2cx \\
x &= \frac{b^2 + c^2 - a^2}{2c} .\\
\end{aligned}
\]

and so
$$
h = \sqrt{b^2 - \left(\frac{b^2 + c^2 - a^2}{2c}\right)^2}.
$$

Now the area of the triangle is
\[
\begin{aligned}
&= \frac12 c \sqrt{b^2 - \left(\frac{b^2 + c^2 - a^2}{2c}\right)^2} \\
&= \sqrt{\frac{b^2c^2}{4} - \frac14\left(\frac{b^2 + c^2 - a^2}{2}\right)^2} \\
&= \sqrt{\left( \frac{bc}{2} \right)^2 - \left(\frac{b^2+c^2-a^2}{4}\right)^2} \\
&= \sqrt{\left(\frac{bc}{2} - \frac{b^2+c^2-a^2}{4}\right)
\left(\frac{bc}{2} + \frac{b^2+c^2-a^2}{4}\right)}\\
&= \sqrt{\left(\frac{a^2-(b^2 - 2bc + c^2)}{4}\right)
\left(\frac{(b^2 + 2bc + c^2) - a^2}{4}\right)}\\
&= \sqrt{\left(\frac{a^2-(b-c)^2}{4}\right)
\left(\frac{(b+c)^2 - a^2}{4}\right)}
\end{aligned}
\]

\[
\begin{aligned}
&= \sqrt{\left(\frac{a-(b-c)}{2}\frac{a+(b-c)}{2}\right)
\left(\frac{(b+c) - a}{2}\frac{(b+c)+a}{2}\right)}\\
&= \sqrt{s(s-a)(s-b)(s-c)}
\end{aligned}
\]

where $s = \frac{a+b+c}{2}$. $\blacksquare$

\end{solution}
\newpage

\begin{problem}{Problem 5}
(a) Consider the following diagram in which the circle has radius 1.
You should assume for this question that you know nothing about the value
of $\pi$.
\begin{center}
\begin{asy}
/* Geogebra to Asymptote conversion, documentation at artofproblemsolving.com/Wiki go to User:Azjps/geogebra */
import graph; size(100);
real labelscalefactor = 0.5; /* changes label-to-point distance */
pen dotstyle = black; /* point style */
real xmin = -2.186177097582655, xmax = 2.4101468649566757, ymin = -1.5794981252782214, ymax = 1.6254094781733246;  /* image dimensions */
pen qqqqtt = rgb(0.,0.,0.2); pen xdxdff = rgb(0.49019607843137253,0.49019607843137253,1.); pen qqzzff = rgb(0.,0.6,1.); pen zzffff = rgb(0.2,1.,0.8);
pair A = (-1.,1.), B = (1.,1.), C = (1.,-1.), D = (-1.,-1.), H_6 = (-1.,0.), H_3 = (1.,0.), H_2 = (0.5,0.8660254037844386), H_1 = (-0.5,0.8660254037844386), H_4 = (0.5,-0.8660254037844386), H_5 = (-0.5,-0.8660254037844386);

filldraw(circle((0.,0.), 1.), xdxdff + opacity(0.15000000596046448), xdxdff);
filldraw(H_1--H_2--H_3--H_4--H_5--H_6--cycle, qqzzff + opacity(0.10000000149011612), qqzzff);
filldraw(A--B--C--D--cycle, zzffff + opacity(0.10000000149011612), zzffff);
/* draw figures */
draw(A--B,  qqqqtt);
draw(B--C );
draw(C--D );
draw(D--A );
draw(H_1--H_2,  qqzzff);
draw(H_2--H_3,  qqzzff);
draw(H_3--H_4,  qqzzff);
draw(H_4--H_5,  qqzzff);
draw(H_5--H_6,  qqzzff);
draw(H_6--H_1,  qqzzff);
draw(H_1--H_2,  qqzzff);
draw(H_2--H_3,  qqzzff);
draw(H_3--H_4,  qqzzff);
draw(H_4--H_5,  qqzzff);
draw(H_5--H_6,  qqzzff);
draw(H_6--H_1,  qqzzff);
draw(A--B,  zzffff);
draw(B--C,  zzffff);
draw(C--D,  zzffff);
draw(D--A,  zzffff);
/* dots and labels */
dot(A,dotstyle);
label("$A$", (-0.985409818030551,1.03808154459437), NE * labelscalefactor);
dot(B,dotstyle);
label("$B$", (1.0172997435459241,1.03808154459437), NE * labelscalefactor);
dot(C,dotstyle);
label("$C$", (1.0172997435459241,-0.9606934205939384), NE * labelscalefactor);
dot(D,dotstyle);
label("$D$", (-0.985409818030551,-0.9606934205939384), NE * labelscalefactor);
dot(H_6,dotstyle);
label("$H_6$", (-0.985409818030551,0.03082486922388395), NE * labelscalefactor);
dot(H_3,dotstyle);
label("$H_3$", (1.0172997435459241,0.03082486922388395), NE * labelscalefactor);
dot(H_2,dotstyle);
label("$H_2$", (0.5176060022488469,0.9043052673967273), NE * labelscalefactor);
dot(H_1,dotstyle);
label("$H_1$", (-0.4857160767334737,0.9043052673967273), NE * labelscalefactor);
dot(H_4,dotstyle);
label("$H_4$", (0.5176060022488469,-0.8269171433962956), NE * labelscalefactor);
dot(H_5,dotstyle);
label("$H_5$", (-0.4857160767334737,-0.8269171433962956), NE * labelscalefactor);
clip((xmin,ymin)--(xmin,ymax)--(xmax,ymax)--(xmax,ymin)--cycle);
/* end of picture */
\end{asy}
\end{center}
i\phantom{i}. Using areas, show that $\pi < 4$. \vspace{0.2cm}

ii. Using perimeters, show that $\pi > 3$. \\

(b) Show that $\frac75 < \sqrt{2} < \frac32$.\vspace{0.2cm}

(c) Show that $\sqrt{2}^{\pi} < \pi^{\sqrt2}$.
\end{problem}

\begin{solution}{Solution}
\textbf{(a) i.} Since the circle is contained within the square, the
area of the square is more than the area of the circle. Thus
$\pi < 2\cdot2=4$. $\square$\\

\textbf{(a) ii.} Let $O$ be the centre of the circle, consider
\begin{center}
\begin{asy}
/* Geogebra to Asymptote conversion, documentation at artofproblemsolving.com/Wiki go to User:Azjps/geogebra */
import graph; size(2.296351105969781cm);
real labelscalefactor = 0.5; /* changes label-to-point distance */
pen dotstyle = black; /* point style */
real xmin = -1.1079994509181452, xmax = 1.1883516550516358, ymin = -0.3471644964790991, ymax = 1.229487297351494;  /* image dimensions */
pen qqzzff = rgb(0.,0.6,1.); pen ffwwqq = rgb(1.,0.4,0.); pen xdxdff = rgb(0.49019607843137253,0.49019607843137253,1.);
pair H_2 = (0.5,0.8660254037844386), H_1 = (-0.5,0.8660254037844386), O = (0.,0.);

filldraw(H_1--H_2--O--cycle, qqzzff + opacity(0.10000000149011612), qqzzff + opacity(0.10000000149011612));
filldraw(arc(O,0.16557936119401293,420.,480.00000000000006)--(0.,0.)--cycle, ffwwqq + opacity(0.10000000149011612),  ffwwqq);
/* draw figures */
draw(H_1--H_2,  qqzzff);
draw(H_1--O,  linetype("4 4") + ffwwqq);
draw(H_2--O,  linetype("4 4") + ffwwqq);
draw(H_1--H_2,  qqzzff);
draw(shift(O)*xscale(1.)*yscale(1.)*arc((0,0),1,60.,120.), xdxdff);
/* dots and labels */
dot(H_2,dotstyle);
label("$H_2$", (0.5073647753502528,0.8828350614311516), NE * labelscalefactor);
dot(H_1,dotstyle);
label("$H_1$", (-0.49211984395306974,0.8828350614311516), NE * labelscalefactor);
dot(O,dotstyle);
label("$O$", (0.00762246569859152,0.01356498760044514), NE * labelscalefactor);
label("$\alpha = 60^\circ$", (-0.006454782460610205,0.1640370760581501), NE * labelscalefactor,ffwwqq);
clip((xmin,ymin)--(xmin,ymax)--(xmax,ymax)--(xmax,ymin)--cycle);
/* end of picture */
\end{asy}
\end{center}
Notice that since $OH_1$ and $OH_2$ are both radii of the same circle,
$\triangle OH_1H_2$ is isosceles which yields $\angle OH_1H_2 =
\angle OH_2H_1 = \frac12(180 - 60)^{\circ} = 60^\circ$. So now we can
conclude that $\triangle OH_1H_2$ is equilateral which gives
$|H_1H_2| = 1$. \vspace{0.2cm}

Since the shortest path from $H_1$ to $H_2$ is the length of the
straight segment $H_1 H_2$ and the arc $\wideparen{OH_1H_2}$ is not equal
to $H_1 H_2$, we have $1 < \frac16(2\pi) = \frac13 \pi \iff 3 < \pi$.
$\square$\\

\textbf{(b)} Squaring the inequality, we have
$$
\frac{49}{25} < 2 < \frac{9}{4}
$$
$$
196 < 200 < 225
$$
which is true. $\square$
\\

\textbf{(c)} We present two solutions,

\textit{Solution 1.}
By (a) i. $\pi < 4$ so $\sqrt{2}^\pi < \sqrt{2}^4 = 4$. By (a) ii. $\pi > 3$
and by (b) $\sqrt2 > \frac75$, and so $\pi^{\sqrt2} > 3^{\frac75}$.
However now observe that
\[
\begin{aligned}
4 &< 3^{\frac75}\\
1024 = 4^5 &< 3^7 = 2187\\
\end{aligned}
\]
which is true.
So we can write $\sqrt{2}^{\pi} < 4 < 3^{\frac75} < \pi^{\sqrt2}$, giving
the desired result.
$\blacksquare$\\

\textit{Solution 2.}
Due to (a) i. $\pi < 4$ and so it suffices to show that
$$
\sqrt{2}^{\pi} < \sqrt{2}^4 = 4 < \pi^{\sqrt2}.
$$
Since $\pi^{\sqrt2} = e^{\sqrt2 \ln \pi}$ we can use the Taylor expansion,
\[
\begin{aligned}
\pi^{\sqrt2} &= \sum_{k=0}^{\infty} \frac{(\sqrt2 \ln \pi)^k}{k!}\\
&> \sum_{k=0}^{\infty} \frac{(\sqrt2)^k}{k!}\\
&> 1 + \sqrt2 + \frac{(\sqrt2)^2}{2} + \frac{(\sqrt2)^3}{6} +
\frac{(\sqrt2)^4}{24}\\
&= \frac{13 + 8\sqrt2}{6}
\end{aligned}
\]
where the second line is because $\pi > 3 > e \implies \ln \pi > 1$. But now note that
\[
\begin{aligned}
\frac{13 + 8\sqrt2}{6} &> 4\\
13 + 8\sqrt2 &> 24\\
8\sqrt2 &> 11\\
128 &> 121
\end{aligned}
\]
and so we have shown that $\pi^{\sqrt2} > 4 > \sqrt2^\pi$, which gives
the desired result. $\blacksquare$

\end{solution}
\newpage

\begin{problem}{Problem 6}
Six non-overlapping squares of side length 2 are contained in a circle. Let
$A$ be the least possible area of a circle allowing this. \vspace{0.2cm}

Show that $A < 12\pi$.
\end{problem}

\begin{solution}{Solution}
Consider the following construction,

\begin{center}
\begin{asy}
/* Geogebra to Asymptote conversion, documentation at artofproblemsolving.com/Wiki go to User:Azjps/geogebra */
import graph; size(320);
real labelscalefactor = 0.5; /* changes label-to-point distance */
pen dps = linewidth(0.7) + fontsize(10); defaultpen(dps); /* default pen style */
pen dotstyle = black; /* point style */
real xmin = -5.93818160963032, xmax = 5.93818160963032, ymin = -1.36921739284875, ymax = 7.84699256425158;  /* image dimensions */
pen ttttff = rgb(0.2,0.2,1.); pen bubtba = rgb(0.7058823529411765,0.7019607843137254,0.7294117647058823); pen rcrcrf = rgb(0.10980392156862745,0.10980392156862745,0.12156862745098039);
pair V_1 = (-2.,0.), V_5 = (-2.,2.), V_7 = (0.,2.), V_2 = (0.,0.), V_9 = (2.,2.), V_3 = (2.,0.), V_8 = (1.,2.), V_6 = (-1.,2.), V_4 = (-3.,2.), gamma  = (0.,2.755901167195395);

filldraw(V_1--V_5--V_4--(-3.,4.)--(-1.,4.)--(-1.,6.)--(1.,6.)--(1.,4.)--(3.,4.)--(3.,2.)--V_9--V_3--cycle, ttttff + opacity(0.10000000149011612),  ttttff);
/* draw grid of horizontal/vertical lines */
pen gridstyle = linewidth(0.7) + bubtba; real gridx = 1., gridy = 1.; /* grid intervals */
for(real i = ceil(xmin/gridx)*gridx; i <= floor(xmax/gridx)*gridx; i += gridx)
draw((i,ymin)--(i,ymax), gridstyle);
for(real i = ceil(ymin/gridy)*gridy; i <= floor(ymax/gridy)*gridy; i += gridy)
draw((xmin,i)--(xmax,i), gridstyle);
/* end grid */

Label laxis; laxis.p = fontsize(10);
xaxis(xmin, xmax,defaultpen+rcrcrf, Ticks(laxis, Step = 1., Size = 2, NoZero),EndArrow(6), above = true);
yaxis(ymin, ymax,defaultpen+rcrcrf, Ticks(laxis, Step = 1., Size = 2, NoZero),EndArrow(6), above = true); /* draws axes; NoZero hides '0' label */
/* draw figures */
draw(V_1--V_5,  ttttff);
draw(V_5--V_4,  ttttff);
draw(V_4--(-3.,4.),  ttttff);
draw((-3.,4.)--(-1.,4.),  ttttff);
draw((-1.,4.)--(-1.,6.),  ttttff);
draw((-1.,6.)--(1.,6.),  ttttff);
draw((1.,6.)--(1.,4.),  ttttff);
draw((1.,4.)--(3.,4.),  ttttff);
draw((3.,4.)--(3.,2.),  ttttff);
draw((3.,2.)--V_9,  ttttff);
draw(V_9--V_3,  ttttff);
draw(V_3--V_1,  ttttff);
draw(V_6--(-1.,4.),  ttttff);
draw((-1.,4.)--(1.,4.),  ttttff);
draw((1.,4.)--V_8,  ttttff);
draw(V_6--V_8,  ttttff);
draw(V_5--V_6,  ttttff);
draw(V_8--V_9,  ttttff);
/* dots and labels */
dot(V_1,dotstyle);
label("$V_1$", (-1.9508523587818016,0.11642297617789425), NE * labelscalefactor);
dot(V_5,dotstyle);
label("$V_5$", (-1.9508523587818016,2.111776483424052), NE * labelscalefactor);
dot(V_7,dotstyle);
label("$V_7$", (0.044501148464345844,2.111776483424052), NE * labelscalefactor);
dot(V_2,dotstyle);
label("$V_2$", (0.044501148464345844,0.09361893609508101), NE * labelscalefactor);
dot(V_9,dotstyle);
label("$V_9$", (2.0512566757519,2.111776483424052), NE * labelscalefactor);
dot(V_3,dotstyle);
label("$V_3$", (2.0512566757519,0.11642297617789425), NE * labelscalefactor);
dot((3.,2.),dotstyle);
label("$V_{10}$", (3.0432324193542706,2.111776483424052), NE * labelscalefactor);
dot((3.,4.),dotstyle);
label("$V_{14}$", (3.0432324193542706,4.118532010711617), NE * labelscalefactor);
dot((1.,4.),dotstyle);
label("$V_{13}$", (1.047878912108123,4.118532010711617), NE * labelscalefactor);
dot(V_8,dotstyle);
label("$V_8$", (1.047878912108123,2.111776483424052), NE * labelscalefactor);
dot(V_6,dotstyle);
label("$V_6$", (-0.9588766151794312,2.111776483424052), NE * labelscalefactor);
dot((-1.,4.),dotstyle);
label("$V_{12}$", (-0.9588766151794312,4.118532010711617), NE * labelscalefactor);
dot(V_4,dotstyle);
label("$V_4$", (-2.954230122425579,2.111776483424052), NE * labelscalefactor);
dot((-3.,4.),dotstyle);
label("$V_{11}$", (-2.954230122425579,4.118532010711617), NE * labelscalefactor);
dot((-1.,6.),dotstyle);
label("$V_{15}$", (-0.9588766151794312,6.113885517957774), NE * labelscalefactor);
dot((1.,6.),dotstyle);
label("$V_{16}$", (1.047878912108123,6.113885517957774), NE * labelscalefactor);
dot(gamma ,dotstyle);
label("$\gamma$", (0.044501148464345844,2.864309806156889), NE * labelscalefactor);
clip((xmin,ymin)--(xmin,ymax)--(xmax,ymax)--(xmax,ymin)--cycle);
/* end of picture */

\end{asy}
\end{center}
where $\gamma = \left(0, 3+\sqrt2 - \frac{\sqrt{11}}{2}\right)$.
\vspace{0.2cm}

The
idea is that we want to place $\gamma$ such that the distance from each
of the vertices is minimised. Since the shape is symmetrical about $x=0$,
we can consider, say only the right side of the polygon. We want to minimise
the distance between $\gamma$ and the outermost vertex that lies at each
$y$ coordinate. For now let $\gamma = (0, k)$. In order to show that
a circle with area less than $12 \pi$ can be constructed we just need to
show that there exists a $k$ such that the maximum of the quantities below
is less than $\sqrt{12}$. We can assume $k > 0$ because obviously we don't
want to place the center below the polygon.
\[
\begin{aligned}
\gamma V_{16} &= \sqrt{1 + (6-k)^2} < \sqrt{12} \implies k > 6-\sqrt{11} \\
\gamma V_{14} &= \sqrt{9 + (4-k)^2} < \sqrt{12} \implies k > 4-\sqrt{3}\\
\gamma V_{10} &= \sqrt{9 + (k-2)^2} < \sqrt{12} \implies k < 2+\sqrt{3}\\
\gamma V_3 &= \sqrt{4 + k^2} < \sqrt{12} \implies k < 2\sqrt{2}
\end{aligned}
\]
Putting together the most restrictive inequalities gives
$6-\sqrt{11} < k < 2\sqrt{2}$. So for any $k$ in this interval a circle
with radius less than $\sqrt{12}$ can be constructed such that it encloses
all six squares. For example picking the midpoint of this interval
$k = \frac12((6-\sqrt{11})+2\sqrt2) =
3 + \sqrt{2} - \sqrt{11}/2$
works. $\blacksquare$
\end{solution}
\newpage

\begin{problem}{Problem 7}
Find all primes $p$ such that $p^2 + 11$ has exactly 6 factors.
\end{problem}

\begin{solution}{Solution}
Since 2 is the only even prime and $2^2 + 11=15$ has exactly 4 factors
we can deduce that $p$ is an odd prime.
By the Fundamental theorem of arithmetic we can write
$$
p^2 + 11 = p_1^{e_1}\cdot p_2^{e_2} \cdots p_n^{e_n}
$$

for some primes $p_1, \ldots, p_n$ and positive integers $e_1, \ldots,
e_n$. Each factor of $p^2 + 11$ is a product of the primes
$p_i^{k_i}$ for integers $1\leq i \leq n$ and $0 \leq k_i \leq e_i$.
Since there are $e_i + 1$ choices for the exponent of $p_i$, we
deduce that $p^2 + 11$ has exactly
$(e_1 + 1)(e_2 + 1) \cdots (e_n + 1)$ factors. So we want to look for
solutions to
$$
(e_1 + 1)(e_2 + 1) \cdots (e_n + 1) = 6.
$$

There are only 2 ways to factor $6 = 3\cdot 2 = 6\cdot 1$. So we must
examine 2 cases.\vspace{0.2cm}

\textit{Case 1.}
If 6 was factored as $6\cdot 1$ then (WLOG) $e_1 = 5$ and all the other
exponents are 0. So our problem reduces to
$$
p^2 + 11 = p_1^5.
$$

But since $p$ is odd, the LHS is even and thus $p_1$ must be even. The only
even prime is 2 so we have $p^2 + 11 = 32$ which gives $p^2 = 21$ which
has no integer solutions. \vspace{0.2cm}

\textit{Case 2.}
If 6 was factored as $3\cdot 2$ then (WLOG) $e_1=2, e_2=1$ and all the
other exponents are 0. The problem becomes
$$
p^2 + 11 = p_1^{2}p_2.
$$

Again since the LHS is even one of those primes must be 2.\vspace{0.2cm}

\textit{Case 2.1.} Suppose that
$p_1 = 2$. Now if $p \neq 3$ then we can use Fermat's little theorem to
show that
$$
p^2 + 11 \equiv 1 + 11 \equiv 0 \equiv 4p_2 \ \text{(mod } 3)
$$

which implies that $p_2 = 3$ since this is the only prime divisible by 3.
So putting that together we have
$p^2 + 11 = 12$ which gives $p = \pm 1$ but this isn't prime.
Now if $p = 3$ then we have $9+11 = 20 = 4p_2$ which gives $p_2 = 5$, and
it can be checked that $p = 3$ does in fact work.\vspace{0.2cm}

\textit{Case 2.2.}
Now suppose that $p_2 = 2$. Proceeding similarly if $p\neq 3$ then
by Fermat's little theorem
$$
p^2 + 11 \equiv 1 + 11 \equiv 0 \equiv 2p_1^2 \ \text{(mod } 3)
$$
which implies that $p_1 = 3$. This gives $p^2 + 11 = 18$ which yields
$p^2 = 7$ which has no integer solutions. Now if $p=3$ then we have
already checked that this leads to a solution, so there is nothing left to
explore.\vspace{0.2cm}

Thus we conclude that the only solution is $p=3$. $\blacksquare$\\

\textit{Comments: This problem looks scary but turns out to be very natural,
in number theory a good number of problems can be solved by just taking
random mods and in this case modulo 3 is enough to solve the problem.}
\end{solution}


\begin{problem}{Problem 8}
\begin{enumerate}[label=(\alph*)]
\item Express 50 in the form $p \pm n^3$ with $p$ prime and $n$ an integer.
\item Find a positive integer that cannot be written in the form
$p \pm n^3$ with $p$ prime and $n$ an integer.
\end{enumerate}
\end{problem}

\begin{solution}{Solution}
\textbf{(a)} $50 = 23 + 3^3$.\\

\textbf{(b)} Because $(-n)^3=-n^3$ we can drop the $\pm$ and just consider
the expression $p + n^3$. Let $b$ be a positive integer such that
$b^3 = p + n^3$. Observe that
$$
p = (b-n)(b^2 + nb + n^2).
$$

Now set $b=8$. Also notice that $b^3 = p + n^3 > n^3 \iff b > n$.
\vspace{0.2cm}

\textit{Case 1.} $b-n = 1$ and $b^2 + nb + n^2 = p$. Now substituting
$b=8$ gives $n = 7$, and so $b^2+nb+n^2 = p = 13^2$. But $p$ is
prime so this is a contradiction. \vspace{0.2cm}

\textit{Case 2.} $b-n = p$ and $b^2 + nb + n^2 = 1$. Substituting $b=8$
gives $n^2 + 8n + 63 = 0$. But since the discriminant
$\Delta = 8^2 - 4(1)(63) = -188 < 0$
there are no real solutions in $n$, contradiction.\vspace{0.2cm}

Thus $8^3 = 512$ is a positive integer that can't be written in the
form $p \pm n^3$. $\blacksquare$ \\

\textit{Comments: Well it feels like there should be some sort of
generalisation here, I'll come back to it later.
}
\end{solution}

\end{document}
